\documentclass[11pt]{article}
\usepackage[a4paper,margin=25mm]{geometry}
\usepackage{amsmath,amssymb,amsthm}
\usepackage[hidelinks]{hyperref}
\emergencystretch=2em
\newtheorem{theorem}{Theorem}
\title{The stated Stern-square normalization tends to zero\newline
\large Disproof of Conjecture 00000000228}
\author{gaochengzhecpu}\date{}
\begin{document}\maketitle
\begin{abstract}
For the actual Stern diatomic sequence, the entire
sequence $s(n^2)/3^{\log_2 n}$ tends to zero as
$n\to\infty$. An exact three-bit recurrence bound
gives the geometric estimate $25(25/27)^m$ on every
block $8^m\leq n<8^{m+1}$. The genuine uniform-index
empirical distributions therefore converge to the
point mass at zero, not to a continuous fractal law.
The proof handles all positive indices, the original
real logarithmic power and actual probability measures.
\end{abstract}

\section*{The original sequence and its adjacent pairs}
Use exactly the defining rules from the source:
\[
 s(0)=0,\quad s(1)=1,\quad s(2n)=s(n),\quad
                  s(2n+1)=s(n)+s(n+1).
\]
Let $v_n=(s(n),s(n+1))$. The rules give
\[
 v_{2n}=L(v_n),\qquad v_{2n+1}=R(v_n),\qquad
 L(a,b)=(a,a+b),\quad R(a,b)=(a+b,b).
\]
These formulas, starting at $v_0=(0,1)$, construct
the actual sequence by repeated division of the
index by two. The Lean project proves all four
original rules, the adjacent-second-coordinate
identity, and uniqueness among sequences satisfying
those rules. It also proves $s(n)>0$ for every $n>0$.

\section*{An exact three-bit growth bound}
If $0\leq a,b\leq B$, every composition of three
maps from $\{L,R\}$ has both coordinates at most
$5B$. Explicitly, writing $v_k=(a,b)$, one obtains
\[
\begin{array}{c|c}
 r&v_{8k+r}\\\hline
 0&(a,3a+b)\\
 1&(3a+b,2a+b)\\
 2&(2a+b,3a+2b)\\
 3&(3a+2b,a+b)\\
 4&(a+b,2a+3b)\\
 5&(2a+3b,a+2b)\\
 6&(a+2b,a+3b)\\
 7&(a+3b,b)
\end{array}
\]
Every displayed coordinate has nonnegative coefficients
with sum at most five. The bound is algebraic and
valid for arbitrary $k$, not inferred from finite testing.

\newpage
\section*{A global bound at every index}
Induction on $j$, dividing an index by eight at each
step, now proves
\[
        0\leq n\leq8^j\quad\Longrightarrow\quad
                 s(n)\leq5^j,
                 \qquad s(n+1)\leq5^j.
\]
At $j=0$ the indices are zero and one, with pairs
$(0,1)$ and $(1,1)$. For the induction step,
$n\leq8^{j+1}$ gives $\lfloor n/8\rfloor\leq8^j$;
three binary steps multiply the pair bound by at
most five. The endpoint is included throughout.

In particular, if $n\leq8^{m+1}$, then
$n^2\leq8^{2(m+1)}$, so
\[
                       s(n^2)\leq5^{2(m+1)}=25^{m+1}.
\]
On the other hand, if $8^m\leq n$, monotonicity
of the ordinary logarithm gives
\[
       \log_2n\geq3m,\qquad 3^{\log_2 n}\geq27^m.
\]
Consequently for every positive index in the block,
\[
 8^m\leq n\leq8^{m+1}
   \quad\Longrightarrow\quad
 0<\frac{s(n^2)}{3^{\log_2n}}
       \leq\frac{25^{m+1}}{27^m}
       =25\left(\frac{25}{27}\right)^m.
\]
For each $n\geq1$, take $m=\lfloor\log_8n\rfloor$.
These blocks cover every positive index, and
$m\to\infty$ with $n$. Since $25/27<1$, we obtain
the full-sequence limit
\[
             a_n:=\frac{s(n^2)}{3^{\log_2n}}
                         \longrightarrow0.
\]
The denominator here is exactly
$3^{\log n/\log2}$, not a rounded exponent or a
replacement normalization. At $n=1$ it is one;
the undefined source index $n=0$ is never sampled.

\section*{The actual empirical distribution}
For sample count $N\geq1$, let
\[
                  \nu_N=\frac1N\sum_{n=1}^{N}\delta_{a_n}.
\]
This is the genuine pushforward of a uniform random
index in $\{1,\ldots,N\}$ through the actual normalized
sequence. All $a_n$ are positive. Thus its CDF
$F_N(x)=\nu_N(( -\infty,x])$ is exactly zero for
every $x\leq0$ and every $N$.

\newpage
\section*{Degeneration and exclusion of a continuous limit}
Fix $x>0$. Since $a_n\to0$, there is $K$ such that
$a_n\leq x$ for every $n\geq K$. At most $K$ of
the first $N$ observations can exceed $x$. Therefore
\[
       0\leq\nu_N((x,\infty))\leq\frac K N\longrightarrow0,
                    \qquad F_N(x)\longrightarrow1.
\]
We have proved the actual CDF limits
\[
 F_N(x)\longrightarrow
       \begin{cases}0,&x<0,\\1,&x>0.\end{cases}
\]
These are precisely the CDF values of the actual
Dirac measure $\delta_0$ at every continuity point
of that CDF. In the usual continuity-point criterion,
the empirical laws converge weakly to $\delta_0$.
No convergence of CDFs to the Dirac value at its
discontinuity zero is asserted; indeed $F_N(0)=0$
for every finite $N$.

There cannot instead be a continuous limiting CDF
$F$. Convergence in law to such a target requires
$F_N(x)\to F(x)$ at every real $x$. It would give
$F(0)=0$ and $F(x)=1$ for every $x>0$. Continuity
along $x_j=1/(j+1)\downarrow0$ would then force
$F(0)=1$, a contradiction.

The necessary continuity-point condition can also
be seen directly from the bounded-continuous-function
definition of weak convergence: approximate the
indicator of $(-\infty,x]$ from above and below
by continuous linear ramps within distance
$\varepsilon$ of $x$, pass to the limit in their
integrals, then let $\varepsilon\downarrow0$ at a
continuity point of the proposed CDF.

\begin{theorem}
For Stern's actual diatomic sequence,
$a_n=s(n^2)/3^{\log_2n}\to0$. Its uniform positive-index
empirical laws have the point mass at zero as their
weak limit and cannot have a continuous limiting CDF.
\end{theorem}
\begin{proof}
The exact three-bit block bound gives the full-sequence
limit. The finite exceptional-index estimate gives
the actual CDF limits, and continuity yields the
displayed contradiction for every continuous target.
\end{proof}
This refutes the source's proposed continuous fractal
distribution under its supplied normalization.
The separate informal spectral-dimension/turning-density
clause need not be assigned an interpretation.
No claim is made about other normalizations of Stern values.

\newpage
\section*{Formalization and reproducibility}
The Lean sequence is constructed by an actual
terminating adjacent-pair recursion. All original
recurrence rules and their uniqueness are proved.
All eight three-step growth cases, the inclusive
power bound, square-index bound and original real
logarithmic denominator comparison are formal theorems.
Mathlib's natural base-eight logarithm supplies blocks
covering every positive integer, and its index tends
to infinity. A geometric squeeze proves the full
normalized limit.

The empirical PMF samples an actual uniform
$\operatorname{Fin}(N+1)$ index and maps it to
$a_{i+1}$. Thus the formal parameter $N\geq0$
represents sample count $N+1$, always using positive
indices. Its canonical probability measure is used
for all subsequent CDF and tail expressions.

The tail probability equals the actual number of
bad indices divided by the sample count. An explicit
finite-set injection bounds that number by $K$.
The exact nonpositive CDF, positive-threshold limits,
equality with Mathlib's actual CDF, and continuity-point
limits to the actual Dirac measure are proved.
A separate theorem rules out every continuous
pointwise CDF limit. The project proves these CDF
statements directly; it does not claim to formalize
the full general equivalence with Mathlib's weak
convergence framework. The necessary condition used
for the disproof is explained on the preceding page.

The supplementary \texttt{verify.py} uses independent
direct-recursion and binary-pair algorithms, checks
all eight block identities at 10,001 base indices,
and checks the square bound at the first 10,000
positive indices, with exact integer/rational arithmetic.
These finite checks are not a Lean computation oracle;
the formal proof supplies the all-index result.

Lean 4.19.0 and Mathlib revision\\
\texttt{c44e0c8ee63ca166450922a373c7409c5d26b00b}
are pinned. From \texttt{lean/}, run
\begin{verbatim}
lake exe cache get
lake build
lake env lean -DwarningAsError=true Main.lean
\end{verbatim}
Fresh-build logs, source and hashes are included.
Printed dependencies contain only standard foundational
axioms. Compile with \texttt{tectonic main.tex}.
The author agent and parent reviewed the work internally;
no external independent review is claimed.

\begin{thebibliography}{9}
\bibitem{source} The Last Math Competition,
Conjecture 00000000228. Exact bilingual source included.
All sequence-specific bounds are proved above and in
the accompanying Lean project.
\end{thebibliography}
\end{document}
